\documentclass[a4paper]{article}
\usepackage[margin=2cm]{geometry}
\usepackage[italian]{babel}
\usepackage[utf8]{inputenc}
\usepackage[T1]{fontenc}
%\usepackage{fancyhdr}
%\usepackage{amsmath, amssymb}
%\usepackage{caption}
\usepackage{graphicx}
%\usepackage{float}
%\pagestyle{fancy}

\begin{document}
\begin{figure}
\centering
\includegraphics[scale=0.07]{../res/images/logo8_crop.png}
\end{figure}

\section*{Verbale incontro 4 Dicembre 2018}

\begin{itemize}
\item Bilanciamento PRO capitolati scartati con quelli del capitolato scelto
nello studio di fattibilità dei vari capitolati. Formalmente i punti a favore 
del capitolato per cui il gruppo intende candidarsi devono essere numericamente
superiori a quelli degli altri capitolati.
\item In merito allo studio di fattibilità, nel paragrafo relativo al dominio
tecnologico vanno inserite le tecnologie obbligatorie, quelle scelte ma non le
tecnologie usabili ma escluse dal gruppo.
\item Possibili ridondanze tra il dominio applicativo e la descrizione generale
del capitolato; in particolare in merito ai capitolati C4 e C5. Si è deciso per
uniformità con gli altri studi di fattibilità di procedere comunque con
l’inserimento di tale paragrafo.
\item Prossimi Compiti / Task
\begin{itemize}
\item Inizio Glossario (Giacomo)
\item Studio di fattibilità C6 (Mattia, Matteo, Samuele)
\item Revisione Studio di Fattibilità Completo (Sara, Francesco)
\item Inizio e definizione struttura norme di progetto (Federico, Paolo)
\end{itemize}
\item Definizione principali / prossime scadenze:
\begin{itemize}
\item 07/12/2018 Terminare Studio di Fattibilità 
\item 11/12/2018 Terminare Prossimi Compiti / Task
\item 09/01/2019 Terminare Lettera di presentazione
\item 10/01/2019 Terminare RR
\end{itemize}
\item Primi contatti con Proponente (Red Babel) e organizzazione primo incontro
o videoconferenza.
\end{itemize}

\end{document}
